\section{Discussion and Limitations}
\label{sec:discussion}

This section interprets the benchmark results of Section~\ref{sec:benchmark} and situates them within the theoretical framework of Sections~\ref{sec:separation}--\ref{sec:bayesian}.

\subsection{What the theory and the benchmark jointly establish}

The structural results (Theorems~\ref{thm:T4}--\ref{thm:T7}) show that, on a compact frequency band and under the collocated-channel condition $CB \neq 0$, the Caputo transfer signature $s^{-\alpha}$ is not reproducible by any finite-dimensional integer-order latent model nor by any strictly proper retarded single-delay model of the class stated in Theorem~\ref{thm:T5}. The finite-horizon analysis (Theorems~\ref{thm:T8}--\ref{thm:T10} and Corollary~\ref{cor:unrestricted}) converts this infinite-bandwidth statement into a quantitative approximation barrier: on any fixed horizon and fixed input, finite exponential mixtures reproduce the fractional convolution output to arbitrary accuracy, so robust discrimination requires an explicit complexity budget on the latent class. The design results of Section~\ref{sec:design} then solve the linear-Gaussian contest exactly: the energy-optimal input is the principal eigenfunction of the noise-weighted inter-model operator (Theorem~\ref{thm:T11}), the robust spectral design needs at most $P+1$ tones (Theorem~\ref{thm:T12}), and the observation schedule is top-$n$ on the separation score (Theorem~\ref{thm:T13}). Sections~\ref{sec:safety} and~\ref{sec:bayesian} add, respectively, the safety certificates (Theorems~\ref{thm:T16}--\ref{thm:T17}) and the exact Bayesian-layer identities (Theorems~\ref{thm:T18}--\ref{thm:T19}) together with a declared theorem/computation boundary.

The benchmark of Section~\ref{sec:benchmark} validates this pipeline in a specific, honest sense: with rivals calibrated to a common stable backbone and a common peak-amplitude budget, an information-criterion (BIC) decision rule attains a macro-averaged accuracy of $0.537$ against a chance level of $0.25$ in the stable memory-shape task. The benchmark does \emph{not} implement the full Bayesian sequential layer of Section~\ref{sec:bayesian}: no posterior model probabilities, expected information gains, evidence estimates, or adaptive iterations are reported. The theoretical and computational layers are therefore consistent but distinct contributions, and the paper keeps them separate.

\subsection{Three levels of the discrimination problem}

A second contribution is methodological: the paper consistently distinguishes three problems that are often conflated.
\begin{enumerate}
\item \emph{Fixed simple hypotheses}: each mechanism is a single fixed predictive distribution. Pairwise KL and eigenvector design (Theorems~\ref{thm:T11}--\ref{thm:T13}) apply directly.
\item \emph{Composite hypotheses}: each mechanism carries nuisance parameters. Robust discrimination requires a maximin criterion over the induced predictive families (Equation~\eqref{eq:bayes_robust}), and positivity of that criterion requires compact parameter sets with the overlap configurations explicitly excluded.
\item \emph{Bayesian model uncertainty}: parameters are integrated out and model identity carries prior mass; this is the level at which the mutual-information identity of Theorem~\ref{thm:T19} is defined and at which adaptive schedules operate.
\end{enumerate}
The benchmark operates at level (1), instantiated through per-model BIC scores; the design theorems are stated at levels (1)--(2); the sequential identities are stated at level (3) and are not executed numerically in this paper.

\subsection{The safety--informativeness trade-off}

The most consequential empirical finding is the trade-off of Section~\ref{sec:benchmark}: at a common amplitude budget, transient designs exhibit much lower Allee-crossing rates but discriminate weakly, while broadband designs discriminate better yet cross the Allee threshold (Figures~\ref{fig:safety-geometry} and~\ref{fig:safety-tradeoff}). At the kernel level, Theorems~\ref{thm:T20}--\ref{thm:T22} provide a protocol-restricted impossibility bound in which approximation error, latent complexity, and Strong-Allee common safety enter the same adversarial hierarchy. At the ecological-response level, Theorem~\ref{thm:T23} and Corollary~\ref{cor:state_testing} work with the exact fractional prey transfer and a finite-state LTI surrogate class; the certified approximation error drops below the kernel-level enclosure from $m=8$ onward. The state-response testing numbers reported in Figure~\ref{fig:state-theorem} use the physical pulse peak cap rather than the broad mass-only safety witness. The latter is retained only as a scope result because the broad response class lacks an independent residue/gain constraint and is not equivalent to the backbone-preserving model family used for ecological interpretation.

Within this regime, the unconstrained eigenvector design of Theorem~\ref{thm:T11} and the feasibility certificate of Theorem~\ref{thm:T16} must be reconciled by a genuine constrained optimization --- not by pointwise saturation of the unconstrained optimum. The projection $u_{\mathrm{safe}}$ of Section~\ref{sec:safety} is a feasible candidate and a lower bound on achievable information, not the constrained optimum; solving the amplitude-capped maximin problem is the natural next step identified by our results.

\subsection{Limitations}

Several limitations must be stated. First, the simulations assume deterministic dynamics corrupted only by additive Gaussian measurement noise; demographic and environmental stochasticity may obscure the spectral signatures exploited here. Second, the finite-horizon barrier is proved for the linearized, fixed-input convolution experiment (Corollary~\ref{cor:unrestricted}), and the certified prey-response theorem of Section~\ref{subsec:state_theorem} is likewise a statement about one linearized input--output channel; extending either to the full nonlinear input--output map requires a Volterra continuous-dependence estimate that we have not developed. Third, the strong-Allee safety certificate is a sufficient rectangle condition with a strict inward margin; it certifies invariance, not optimality of the admissible set. Fourth, the latent-dimension cap $m \le 5$ is a declared modeling budget with sensitivity analysis, not a biological law. Fifth, the decision rule is BIC --- an information criterion --- and the Bayesian layer of Section~\ref{sec:bayesian} is specified, not executed.

The paper does not claim novelty for safe information limits, safe active model discrimination, or finite Prony/rational approximation in isolation. Contemporary work already studies information limits induced by safety and safe model-discrimination policies under state/input constraints~\cite{SAFE26A,SAFE26B}, while pole-aware Mittag--Leffler approximation and the finite-Prony versus fractional branch-cut distinction are established in adjacent literatures~\cite{ML01,PRONY26}. Relative to optimal experimental design for dynamical systems~\cite{O01,O05} and fractional-order identification~\cite{I02,F02}, our narrower contribution is the joint construction: a certified strong-Allee ecological backbone; explicit fractional-versus-latent finite-horizon approximation hierarchy; protocol-restricted safety interface; interval-certified prey-response finite-state enclosures; and testing lower bounds that quantify how increasing latent complexity erodes identifiable fractional-memory evidence.

\subsection{Future work}

The findings point to four concrete extensions: (i) the safety-constrained optimal design problem itself --- amplitude-capped eigenvector/maximin optimization inside a rigorously verified safe set; (ii) the full Bayesian sequential layer --- evidence estimation, expected information gain via the identity of Theorem~\ref{thm:T19}, adaptive scheduling, and simulation-based calibration; (iii) heterogeneous memory classes --- distributed-order derivatives and multiple delay scales under the same complexity-budget discipline; and (iv) experimental validation in microbial microcosms following the protocols of Section~\ref{sec:prospect}.
